\section{Dispositivos Virtuais}

Dispositivos virtuais representam, para o \gls{so} convidado, dispositivos de hardware que podem ser utilizados para comunicação com a plataforma,
como temporizadores e dispositivos de \gls{io}. Eles abstraem os mecanismos de hardware utilizados pelo hipervisor, permitindo que o convidado interaja com esses
recursos sem acessar diretamente os dispositivos físicos.

\subsection{Comunicação: Convidado para Hardware}

A comunicação do \gls{so} convidado com o hardware é intermediada pelo hipervisor. As requisições do convidado podem chegar ao hipervisor tanto por meio da interceptação de operações
de acesso ao hardware quanto por chamadas direcionadas diretamente aos dispositivos virtuais. Em ambos os casos, o hipervisor encaminha a requisição ao objeto da \gls{vm} correspondente, que
identifica o dispositivo virtual responsável pela operação. O dispositivo virtual, por sua vez, traduz e encaminha a requisição ao mediador de hardware correspondente,
conforme ilustrado na Figura~\ref{fig:diagram-hosted-to-hardware}.

\begin{figure}[H]
\centering
\includegraphics[width=1\textwidth]{images/VirtualDeviceToHardware.pdf}
\caption{Diagrama do Fluxo de Comunicação: Convidado para Hardware}\label{fig:diagram-hosted-to-hardware}
\end{figure}

\subsection{Comunicação: Hardware para Convidado}

O fluxo de notificação de um evento de hardware para o \gls{so} convidado é ilustrado na Figura~\ref{fig:diagram-hardware-to-hosted}.
Esse mecanismo baseia-se no padrão de projeto \textit{Observer}, no qual cada mediador de hardware atua como um elemento observável e cada dispositivo virtual funciona como um observador.

Quando ocorre um evento no mediador de hardware, este notifica seus observadores cadastrados, propagando a sinalização até que o dispositivo virtual correspondente
processe a informação e a entregue ao domínio do SO convidado.

\begin{figure}[H]
\centering
\includegraphics[width=1\textwidth]{images/HardwareToVirtualDevice.pdf}
\caption{Diagrama do Fluxo de Comunicação: Hardware para Convidado}\label{fig:diagram-hardware-to-hosted}
\end{figure}
